\usepackage{needspace}
\usepackage{caption}
\usepackage{chngcntr}
\usepackage{xcolor}
\usepackage{framed}
\usepackage{fvextra}

% A4技術書向けの本文段落設定。
% 段落頭は字下げせず、段落間は0.3emだけ空ける。
\setlength{\parindent}{0pt}
\setlength{\parskip}{0.3em}

% PDFに残るコードフェンスは薄いグレー背景と外枠を付ける。
% powershell, text, yaml, markdown, json, html, css, javascript, python, latex, r を含む。
% mermaidは既存処理で図へ変換されるため、通常はコードブロックとしてPDFへ残らない。
\definecolor{bookcodebackground}{RGB}{245,245,245}
\definecolor{bookcodeborder}{RGB}{153,153,153}
% Verbatim行はベースライン基準のため、見た目に合わせて上下を個別補正する。
% 上側は従来の3pt補正を維持し、下側だけ0.75pt詰めて2.25ptとする。
\newlength{\bookcodetopadjust}
\newlength{\bookcodebottomadjust}
\setlength{\bookcodetopadjust}{3pt}
\setlength{\bookcodebottomadjust}{2.25pt}
\newenvironment{bookcodeblock}{%
  \def\FrameCommand{%
    \fboxsep=6pt\fboxrule=0.6pt\fcolorbox{bookcodeborder}{bookcodebackground}%
  }%
  \MakeFramed{\advance\hsize-\width\FrameRestore}%
  \vspace*{-\bookcodetopadjust}%
}{%
  \vspace*{\bookcodebottomadjust}%
  \endMakeFramed
}

% 図番号は章をまたいで連番とし、キャプションは図の下へ置く。
% 文字サイズは本文と同じ大きさにする。
\counterwithout{figure}{chapter}
\renewcommand{\thefigure}{\arabic{figure}}
\captionsetup[figure]{
  name=図,
  position=bottom,
  font=normalsize,
  labelfont=bf,
  labelsep=space
}

% 表番号も章をまたいで連番とし、キャプションは表の上へ置く。
% 「表1」のように名称と番号の間を空けずに表示する。
\counterwithout{table}{chapter}
\renewcommand{\thetable}{\arabic{table}}
\DeclareCaptionLabelFormat{booktable}{#1#2}
\captionsetup[table]{
  name=表,
  position=top,
  font=normalsize,
  labelfont=bf,
  labelformat=booktable,
  labelsep=space
}

% 3列以上の表で使用するデータ行区切りの色。
% 紙面を重くしないよう、かなり淡いグレーとする。
\definecolor{booktablerule}{RGB}{210,210,210}

% Markdownの > による引用を、コードとは区別した淡い引用ボックスにする。
\definecolor{bookquotebackground}{RGB}{250,250,250}
\definecolor{bookquoteborder}{RGB}{180,180,180}
\newenvironment{bookquoteblock}{%
  \def\FrameCommand{%
    \fboxsep=8pt\fboxrule=0.6pt\fcolorbox{bookquoteborder}{bookquotebackground}%
  }%
  \MakeFramed{\advance\hsize-\width\FrameRestore}%
}{%
  \endMakeFramed
}

\makeatletter

% 目次では柱とページ番号を表示しない
\newif\ifbooktoc
\booktocfalse

\def\ps@bookfront{%
  \let\@mkboth\@gobbletwo
  \let\@oddhead\@empty
  \let\@evenhead\@empty
  \let\@oddfoot\@empty
  \let\@evenfoot\@empty
}

% 章開始ページ専用の柱
% 奇数ページは見出しを左、ページ番号を右に置く
% 偶数ページはページ番号を左、見出しを右に置く
\newcommand{\bookchaptertitle}{}
\def\ps@bookchapter{%
  \let\@mkboth\@gobbletwo
  \def\@oddhead{%
    \vbox{%
      \hbox to\textwidth{\bookchaptertitle\hfil\thepage}%
      \vskip 2pt
      \hrule
    }%
  }%
  \def\@evenhead{%
    \vbox{%
      \hbox to\textwidth{\thepage\hfil\bookchaptertitle}%
      \vskip 2pt
      \hrule
    }%
  }%
  \let\@oddfoot\@empty
  \let\@evenfoot\@empty
}

% タイトルページは書名、副題、著者名だけを表示する。
% 発行日と著作権表示は奥付へ集約する。

\renewcommand{\@makechapterhead}[1]{%
  \markboth{第\thechapter 章\space #1}{第\thechapter 章\space #1}%
  \gdef\bookchaptertitle{第\thechapter 章\space #1}%
  \thispagestyle{bookchapter}%
  \vspace*{-10pt}%
  {\parindent \z@ \raggedright
    \normalfont\sffamily\bfseries\Large
    第\thechapter 章\quad #1\par\nobreak
    \vskip 24pt
  }%
}

\renewcommand{\@makeschapterhead}[1]{%
  \ifbooktoc
    \thispagestyle{bookfront}%
  \else
    \markboth{#1}{#1}%
    \markright{#1}%
    \gdef\bookchaptertitle{#1}%
    \thispagestyle{bookchapter}%
  \fi
  \vspace*{-10pt}%
  {\parindent \z@ \raggedright
    \normalfont\sffamily\bfseries\Large #1\par\nobreak
    \vskip 24pt
  }%
}

\AtBeginDocument{%
  % MarkdownのBlockQuote（>）を引用ボックスへ変換する。
  \renewenvironment{quote}{\begin{bookquoteblock}}{\end{bookquoteblock}}%
  \renewenvironment{quotation}{\begin{bookquoteblock}}{\end{bookquoteblock}}%

  % コードブロックの長い行は右端で自動改行する。
  % breakanywhere=true により、空白のない長い文字列も枠内で折り返す。
  % listparameters で Verbatim が持つ上下の余白を消し、
  % 外枠側の fboxsep=6pt だけを残して天地方向の余白を揃える。
  \RecustomVerbatimEnvironment{verbatim}{Verbatim}{%
    breaklines=true,
    breakanywhere=true,
    breaksymbolleft={},
    breaksymbolright={},
    listparameters={\setlength{\topsep}{0pt}\setlength{\partopsep}{0pt}\setlength{\parsep}{0pt}},
    frame=single,
    framesep=6pt,
    framerule=0.6pt,
    rulecolor=\color{bookcodeborder},
    fillcolor=\color{bookcodebackground}
  }%
  % Pandocのシンタックスハイライト付きコードも同じ折り返し規則を使う。
  % commandchars はPandocの色付け用マクロを有効にするため維持する。
  \@ifundefined{Highlighting}{}{%
    \RecustomVerbatimEnvironment{Highlighting}{Verbatim}{%
      commandchars=\\\{\},
      breaklines=true,
      breakanywhere=true,
      breaksymbolleft={},
      breaksymbolright={},
      listparameters={\setlength{\topsep}{0pt}\setlength{\partopsep}{0pt}\setlength{\parsep}{0pt}}
    }%
  }%
  % シンタックスハイライト付きコードフェンスの外観を統一する。
  \@ifundefined{Shaded}{}{%
    \renewenvironment{Shaded}{\begin{bookcodeblock}}{\end{bookcodeblock}}%
  }%
  % 章開始ページのplainを、その時点のheadingsへ動的に接続する
  \def\ps@plain{\ps@headings}%
  % 目次だけは柱を消し、終了後に通常のheadingsへ戻す
  \let\bookoriginaltableofcontents\tableofcontents
  \renewcommand{\tableofcontents}{%
    \clearpage
    \booktoctrue
    \pagestyle{bookfront}%
    \thispagestyle{bookfront}%
    \let\ps@plain\ps@bookfront
    \bookoriginaltableofcontents
    \clearpage
    \booktocfalse
    \def\ps@plain{\ps@headings}%
    \pagestyle{headings}%
  }%
  % 図キャプションの名称、位置、文字サイズはcaptionパッケージで統一する
  \hypersetup{pdfpagemode=UseNone,pdfpagelayout=SinglePage}
  \pdfextension catalog{/PageMode /UseNone}
  \let\originalmaketitle\maketitle
  \renewcommand{\maketitle}{%
    \begingroup
      % dateメタデータは奥付で使用するが、タイトルページには表示しない
      \let\@date\@empty
      {\sffamily\originalmaketitle}%
    \endgroup
    \clearpage
  }%
}

\makeatother
